\documentclass[12pt,a4paper]{article}

% Estilo compartido (paquetes, colores, hyperref, listings, AI Supervisor)
%% ============================================================
%% estilo_3dosim.tex — Estilo compartido para documentos 3Dosim
%% ============================================================
%% Incluir con: %% ============================================================
%% estilo_3dosim.tex — Estilo compartido para documentos 3Dosim
%% ============================================================
%% Incluir con: %% ============================================================
%% estilo_3dosim.tex — Estilo compartido para documentos 3Dosim
%% ============================================================
%% Incluir con: \input{../estilo_3dosim}
%% ============================================================

%% ============================================================
%% Paquetes base
%% ============================================================
\usepackage[utf8]{inputenc}
\usepackage[T1]{fontenc}
\usepackage[spanish]{babel}
\usepackage{amsmath,amssymb,amsfonts}
\usepackage{graphicx}
\usepackage{xcolor}
\usepackage{hyperref}
\usepackage{geometry}
\usepackage{fancyhdr}
\usepackage{titlesec}
\usepackage{float}
\usepackage{caption}
\usepackage{booktabs}
\usepackage{longtable}
\usepackage{array}
\usepackage[most]{tcolorbox}
\usepackage{listings}

\geometry{margin=2.5cm}
\setlength{\headheight}{13.6pt}

%% ============================================================
%% PALETA DE COLORES 3Dosim
%% ============================================================
\definecolor{cPrimary}{RGB}{30,64,124}      % azul oscuro titulos
\definecolor{cSecondary}{RGB}{70,130,180}   % azul claro subtitulos
\definecolor{cAccent}{RGB}{220,50,50}        % rojo AI Supervisor / alertas
\definecolor{cBg}{RGB}{245,245,245}          % gris muy claro fondo
\definecolor{cDarkBg}{RGB}{230,230,230}      % gris fondo tablas

%% Colores para listings (sintaxis)
\definecolor{codegreen}{RGB}{28,172,0}
\definecolor{codegray}{RGB}{128,128,128}
\definecolor{codepurple}{RGB}{170,50,170}
\definecolor{codeblue}{RGB}{21,100,210}

%% ============================================================
%% FORMATO DE SECCIONES CON COLORES
%% ============================================================
\titleformat*{\section}{\normalfont\Large\bfseries\color{cPrimary}}
\titleformat*{\subsection}{\normalfont\large\bfseries\color{cSecondary}}
\titleformat*{\subsubsection}{\normalfont\normalsize\bfseries\color{cPrimary}}

%% ============================================================
%% CONFIGURACION hyperref — enlaces con color y metadatos PDF
%% ============================================================
\hypersetup{
  colorlinks=true,
  linkcolor=cPrimary,
  citecolor=cSecondary,
  urlcolor=cSecondary,
  bookmarksopen=true,
  pdfcreator={3Dosim LaTeX},
  pdfsubject={Dosimetría Hepática Y-90}
}

%% ============================================================
%% CONFIGURACION fancyhdr — encabezado/pie base
%% ============================================================
\pagestyle{fancy}
\fancyhf{}
\fancyfoot[C]{\thepage}
\renewcommand{\headrulewidth}{0.4pt}

%% ============================================================
%% Macro para notacion: \vardef{$A_{vox}$}{Descripcion}
%% ============================================================
\newcommand{\vardef}[2]{\noindent\textbf{#1}: #2\\}
\setlength{\parskip}{0.5em}

%% ============================================================
%% CAJA COLOREADA PARA AI SUPERVISOR
%% Uso: \begin{tcolorbox}[aiSupervisor] ... \end{tcolorbox}
%% ============================================================
\tcbset{
  aiSupervisor/.style = {
    colback=cBg,
    colframe=cAccent,
    boxrule=0.7pt,
    arc=3pt,
    left=4mm,
    right=4mm,
    top=3mm,
    bottom=3mm,
    fonttitle=\bfseries\large,
    title=Control de Calidad (AI Supervisor),
    coltitle=white,
    colbacktitle=cAccent,
    attach boxed title to top left={yshift=-3mm,xshift=5mm},
    boxed title style={colback=cAccent,arc=2pt}
  }
}

%% ============================================================
%% CONFIGURACION listings — reemplazo de verbatim con resaltado
%% ============================================================
\lstset{
  language=Python,
  basicstyle=\small\ttfamily,
  keywordstyle=\color{codeblue}\bfseries,
  commentstyle=\color{codegray}\itshape,
  stringstyle=\color{codepurple},
  numbers=left,
  numberstyle=\tiny\color{codegray},
  stepnumber=1,
  numbersep=5pt,
  backgroundcolor=\color{cBg},
  frame=single,
  rulecolor=\color{cDarkBg},
  breaklines=true,
  breakatwhitespace=false,
  tabsize=2,
  showstringspaces=false,
  captionpos=b,
  literate={á}{{\'a}}1 {é}{{\'e}}1 {í}{{\'i}}1 {ó}{{\'o}}1 {ú}{{\'u}}1
           {Á}{{\'A}}1 {É}{{\'E}}1 {Í}{{\'I}}1 {Ó}{{\'O}}1 {Ú}{{\'U}}1
           {ñ}{{\~n}}1 {Ñ}{{\~N}}1 {ü}{{\"u}}1
}

%% ============================================================

%% ============================================================
%% Paquetes base
%% ============================================================
\usepackage[utf8]{inputenc}
\usepackage[T1]{fontenc}
\usepackage[spanish]{babel}
\usepackage{amsmath,amssymb,amsfonts}
\usepackage{graphicx}
\usepackage{xcolor}
\usepackage{hyperref}
\usepackage{geometry}
\usepackage{fancyhdr}
\usepackage{titlesec}
\usepackage{float}
\usepackage{caption}
\usepackage{booktabs}
\usepackage{longtable}
\usepackage{array}
\usepackage[most]{tcolorbox}
\usepackage{listings}

\geometry{margin=2.5cm}
\setlength{\headheight}{13.6pt}

%% ============================================================
%% PALETA DE COLORES 3Dosim
%% ============================================================
\definecolor{cPrimary}{RGB}{30,64,124}      % azul oscuro titulos
\definecolor{cSecondary}{RGB}{70,130,180}   % azul claro subtitulos
\definecolor{cAccent}{RGB}{220,50,50}        % rojo AI Supervisor / alertas
\definecolor{cBg}{RGB}{245,245,245}          % gris muy claro fondo
\definecolor{cDarkBg}{RGB}{230,230,230}      % gris fondo tablas

%% Colores para listings (sintaxis)
\definecolor{codegreen}{RGB}{28,172,0}
\definecolor{codegray}{RGB}{128,128,128}
\definecolor{codepurple}{RGB}{170,50,170}
\definecolor{codeblue}{RGB}{21,100,210}

%% ============================================================
%% FORMATO DE SECCIONES CON COLORES
%% ============================================================
\titleformat*{\section}{\normalfont\Large\bfseries\color{cPrimary}}
\titleformat*{\subsection}{\normalfont\large\bfseries\color{cSecondary}}
\titleformat*{\subsubsection}{\normalfont\normalsize\bfseries\color{cPrimary}}

%% ============================================================
%% CONFIGURACION hyperref — enlaces con color y metadatos PDF
%% ============================================================
\hypersetup{
  colorlinks=true,
  linkcolor=cPrimary,
  citecolor=cSecondary,
  urlcolor=cSecondary,
  bookmarksopen=true,
  pdfcreator={3Dosim LaTeX},
  pdfsubject={Dosimetría Hepática Y-90}
}

%% ============================================================
%% CONFIGURACION fancyhdr — encabezado/pie base
%% ============================================================
\pagestyle{fancy}
\fancyhf{}
\fancyfoot[C]{\thepage}
\renewcommand{\headrulewidth}{0.4pt}

%% ============================================================
%% Macro para notacion: \vardef{$A_{vox}$}{Descripcion}
%% ============================================================
\newcommand{\vardef}[2]{\noindent\textbf{#1}: #2\\}
\setlength{\parskip}{0.5em}

%% ============================================================
%% CAJA COLOREADA PARA AI SUPERVISOR
%% Uso: \begin{tcolorbox}[aiSupervisor] ... \end{tcolorbox}
%% ============================================================
\tcbset{
  aiSupervisor/.style = {
    colback=cBg,
    colframe=cAccent,
    boxrule=0.7pt,
    arc=3pt,
    left=4mm,
    right=4mm,
    top=3mm,
    bottom=3mm,
    fonttitle=\bfseries\large,
    title=Control de Calidad (AI Supervisor),
    coltitle=white,
    colbacktitle=cAccent,
    attach boxed title to top left={yshift=-3mm,xshift=5mm},
    boxed title style={colback=cAccent,arc=2pt}
  }
}

%% ============================================================
%% CONFIGURACION listings — reemplazo de verbatim con resaltado
%% ============================================================
\lstset{
  language=Python,
  basicstyle=\small\ttfamily,
  keywordstyle=\color{codeblue}\bfseries,
  commentstyle=\color{codegray}\itshape,
  stringstyle=\color{codepurple},
  numbers=left,
  numberstyle=\tiny\color{codegray},
  stepnumber=1,
  numbersep=5pt,
  backgroundcolor=\color{cBg},
  frame=single,
  rulecolor=\color{cDarkBg},
  breaklines=true,
  breakatwhitespace=false,
  tabsize=2,
  showstringspaces=false,
  captionpos=b,
  literate={á}{{\'a}}1 {é}{{\'e}}1 {í}{{\'i}}1 {ó}{{\'o}}1 {ú}{{\'u}}1
           {Á}{{\'A}}1 {É}{{\'E}}1 {Í}{{\'I}}1 {Ó}{{\'O}}1 {Ú}{{\'U}}1
           {ñ}{{\~n}}1 {Ñ}{{\~N}}1 {ü}{{\"u}}1
}

%% ============================================================

%% ============================================================
%% Paquetes base
%% ============================================================
\usepackage[utf8]{inputenc}
\usepackage[T1]{fontenc}
\usepackage[spanish]{babel}
\usepackage{amsmath,amssymb,amsfonts}
\usepackage{graphicx}
\usepackage{xcolor}
\usepackage{hyperref}
\usepackage{geometry}
\usepackage{fancyhdr}
\usepackage{titlesec}
\usepackage{float}
\usepackage{caption}
\usepackage{booktabs}
\usepackage{longtable}
\usepackage{array}
\usepackage[most]{tcolorbox}
\usepackage{listings}

\geometry{margin=2.5cm}
\setlength{\headheight}{13.6pt}

%% ============================================================
%% PALETA DE COLORES 3Dosim
%% ============================================================
\definecolor{cPrimary}{RGB}{30,64,124}      % azul oscuro titulos
\definecolor{cSecondary}{RGB}{70,130,180}   % azul claro subtitulos
\definecolor{cAccent}{RGB}{220,50,50}        % rojo AI Supervisor / alertas
\definecolor{cBg}{RGB}{245,245,245}          % gris muy claro fondo
\definecolor{cDarkBg}{RGB}{230,230,230}      % gris fondo tablas

%% Colores para listings (sintaxis)
\definecolor{codegreen}{RGB}{28,172,0}
\definecolor{codegray}{RGB}{128,128,128}
\definecolor{codepurple}{RGB}{170,50,170}
\definecolor{codeblue}{RGB}{21,100,210}

%% ============================================================
%% FORMATO DE SECCIONES CON COLORES
%% ============================================================
\titleformat*{\section}{\normalfont\Large\bfseries\color{cPrimary}}
\titleformat*{\subsection}{\normalfont\large\bfseries\color{cSecondary}}
\titleformat*{\subsubsection}{\normalfont\normalsize\bfseries\color{cPrimary}}

%% ============================================================
%% CONFIGURACION hyperref — enlaces con color y metadatos PDF
%% ============================================================
\hypersetup{
  colorlinks=true,
  linkcolor=cPrimary,
  citecolor=cSecondary,
  urlcolor=cSecondary,
  bookmarksopen=true,
  pdfcreator={3Dosim LaTeX},
  pdfsubject={Dosimetría Hepática Y-90}
}

%% ============================================================
%% CONFIGURACION fancyhdr — encabezado/pie base
%% ============================================================
\pagestyle{fancy}
\fancyhf{}
\fancyfoot[C]{\thepage}
\renewcommand{\headrulewidth}{0.4pt}

%% ============================================================
%% Macro para notacion: \vardef{$A_{vox}$}{Descripcion}
%% ============================================================
\newcommand{\vardef}[2]{\noindent\textbf{#1}: #2\\}
\setlength{\parskip}{0.5em}

%% ============================================================
%% CAJA COLOREADA PARA AI SUPERVISOR
%% Uso: \begin{tcolorbox}[aiSupervisor] ... \end{tcolorbox}
%% ============================================================
\tcbset{
  aiSupervisor/.style = {
    colback=cBg,
    colframe=cAccent,
    boxrule=0.7pt,
    arc=3pt,
    left=4mm,
    right=4mm,
    top=3mm,
    bottom=3mm,
    fonttitle=\bfseries\large,
    title=Control de Calidad (AI Supervisor),
    coltitle=white,
    colbacktitle=cAccent,
    attach boxed title to top left={yshift=-3mm,xshift=5mm},
    boxed title style={colback=cAccent,arc=2pt}
  }
}

%% ============================================================
%% CONFIGURACION listings — reemplazo de verbatim con resaltado
%% ============================================================
\lstset{
  language=Python,
  basicstyle=\small\ttfamily,
  keywordstyle=\color{codeblue}\bfseries,
  commentstyle=\color{codegray}\itshape,
  stringstyle=\color{codepurple},
  numbers=left,
  numberstyle=\tiny\color{codegray},
  stepnumber=1,
  numbersep=5pt,
  backgroundcolor=\color{cBg},
  frame=single,
  rulecolor=\color{cDarkBg},
  breaklines=true,
  breakatwhitespace=false,
  tabsize=2,
  showstringspaces=false,
  captionpos=b,
  literate={á}{{\'a}}1 {é}{{\'e}}1 {í}{{\'i}}1 {ó}{{\'o}}1 {ú}{{\'u}}1
           {Á}{{\'A}}1 {É}{{\'E}}1 {Í}{{\'I}}1 {Ó}{{\'O}}1 {Ú}{{\'U}}1
           {ñ}{{\~n}}1 {Ñ}{{\~N}}1 {ü}{{\"u}}1
}


% Encabezado especifico del modulo
\fancyhead[L]{\small 3Dosim — Cálculo de Dosis}
\fancyhead[R]{\small MCTALL y Kernel}

% Portada
\title{\textbf{Cálculo de Dosis en 3Dosim:} \\
       Método MCTALL (Monte Carlo) y \\
       Método Kernel (Convolución)}
\author{Proyecto 3Dosim — Dosimetría Hepática con $^{90}$Y}
\date{\today}

\begin{document}

\maketitle
\thispagestyle{empty}
\newpage

% ============================================================
% Tabla de contenidos
% ============================================================
\tableofcontents
\newpage

% ============================================================
% INTRODUCCIÓN
% ============================================================
\section{Introducción}

El proyecto \textbf{3Dosim} implementa dos metodologías complementarias para el
cálculo de dosis absorbida en tratamientos de radioembolización hepática con
$^{90}$Y (Ytrio-90):

\begin{enumerate}
    \item \textbf{Método MCTALL (Monte Carlo):} Simulación completa del
        transporte de radiación mediante MCNP (Monte Carlo N-Particle). La
        dosis se obtiene directamente del \emph{mesh tally} FMESH4, que
        registra la energía depositada por partícula fuente en cada voxel.
    \item \textbf{Método Kernel (Convolución):} Aproximación analítica que
        convoluciona la imagen PET de actividad con un \emph{kernel de dosis}
        precalculado por MCNP para una fuente puntual en medio infinito
        uniforme.
\end{enumerate}

Ambos métodos convergen al mismo resultado físico cuando la estadística de
Monte Carlo es suficiente, pero difieren en costo computacional y
aplicabilidad clínica.

% ============================================================
% FUNDAMENTOS FÍSICOS
% ============================================================
\section{Fundamentos Físicos}

\subsection{El isótopo $^{90}$Y}

El $^{90}$Y es un emisor $\beta^-$ puro con las siguientes características:

\begin{itemize}
    \item Vida media: $T_{1/2} = 64.1$ horas
    \item Constante de desintegración: $\lambda = \frac{\ln 2}{T_{1/2}}$
    \item Vida media (mean lifetime): $\tau = \frac{1}{\lambda} = \frac{T_{1/2}}{\ln 2} \approx 332\,916$ s $\approx 92.43$ h
    \item Energía máxima $\beta^-$: $E_{\max} = 2.28$ MeV
    \item Energía media $\beta^-$: $\bar{E}_\beta = 0.9336$ MeV
    \item Alcance máximo en tejido: $\sim 11$ mm
    \item Alcance medio: $\sim 2.5$ mm
\end{itemize}

La corta vida media y el alcance limitado hacen del $^{90}$Y un isótopo
ideal para radioembolización, ya que la energía se deposita localmente en
el tejido tumoral.

\subsection{Dosis absorbida}

La dosis absorbida $D$ se define como la energía depositada por unidad de
masa:

\begin{equation}
    D = \frac{dE}{dm} \quad [\text{J/kg} = \text{Gy}]
\end{equation}

En dosimetría interna, la dosis en un voxel depende de:
\begin{itemize}
    \item La actividad acumulada en el voxel (fuente)
    \item La energía liberada por desintegración ($\bar{E}_\beta$)
    \item La fracción de energía absorbida localmente (en $^{90}$Y,
        esencialmente 1 para $\beta^-$)
    \item La masa del voxel (densidad $\times$ volumen)
\end{itemize}

% ============================================================
% MÉTODO MCTALL
% ============================================================
\section{Método MCTALL: Monte Carlo con MCNP}

\subsection{Descripción general}

El método MCTALL utiliza el código de Monte Carlo \textbf{MCNP} (Monte Carlo
N-Particle) para simular el transporte de las partículas $\beta^-$ emitidas
por el $^{90}$Y a través de una geometría voxelizada del paciente
(\emph{phantom}).

El flujo de trabajo es:

\begin{enumerate}
    \item Se construye un \emph{phantom} 3D voxelizado a partir de la
        segmentación del CT del paciente, donde cada voxel tiene asignado
        un material MCNP (tejido blando, pulmón, hueso, hígado, tumor, etc.)
        con su densidad y composición elemental.
    \item Se define la fuente (SDEF) a partir de la imagen PET, que
        proporciona la distribución de actividad.
    \item Se define un \emph{mesh tally} FMESH4 que cubre todo el volumen
        de interés. El tally registra la energía depositada por partícula
        fuente en cada voxel.
    \item Se ejecuta MCNP con NPS (number of particle histories) suficiente
        para obtener estadística aceptable.
    \item Se lee el archivo de salida \textbf{MCTAL} y se extrae la dosis
        3D.
    \item Se convierte de MeV/cm$^3$/partícula a Gy.
\end{enumerate}

\subsection{Estructura del archivo MCTAL}

El archivo MCTAL es la salida ASCII de MCNP que contiene los resultados de
los tallies. Para el FMESH4, la estructura relevante es:

\begin{lstlisting}
tally    1   -1   -1
 0 0 1 0 0 0 0 0 0 ...
 f <total> <type> <nx> <ny> <nz>
 <x boundaries nx+1>
 <y boundaries ny+1>
 <z boundaries nz+1>
 t        1
 vals
 v1 e1 v2 e2 v3 e3 v4 e4
 v5 e5 v6 e6 v7 e7 v8 e8
 ...
\end{lstlisting}

Donde:
\begin{itemize}
    \item \texttt{nx, ny, nz}: dimensiones del mesh tally
    \item \texttt{v$_i$}: valor de dosis en MeV/cm$^3$ por partícula fuente
    \item \texttt{e$_i$}: incertidumbre relativa asociada
\end{itemize}

\subsection{Parsing del MCTAL}

El parsing sigue el algoritmo de \texttt{f\_cargo\_mctall.m} (MATLAB) y
\texttt{MCTALParser} (Python):

\begin{enumerate}
    \item Buscar la línea \texttt{tally    1} en el archivo.
    \item Leer la línea \texttt{f} para obtener dimensiones $(nx, ny, nz)$.
    \item Saltar las líneas de boundaries y el marcador \texttt{vals}.
    \item Leer los pares (valor, error) como floats: el archivo los
        almacena en grupos de 8 columnas.
    \item Reconstruir el array 3D:
        \begin{itemize}
            \item Los valores pares ($0,2,4,\dots$) son dosis.
            \item Los valores impares ($1,3,5,\dots$) son errores.
            \item Reshape column-major a $[ny, nx, nz]$ (convención
                Fortran/MATLAB).
            \item Transponer cada slice: $[:,:,i] \to [:,:,i]^T$, obteniendo
                $[nx, ny, nz]$.
        \end{itemize}
\end{enumerate}

\subsection{Conversión de MeV/cm$^3$ a Gy}

La conversión de la dosis cruda del MCTAL ($D_{\text{raw}}$ en
MeV/cm$^3$/partícula) a dosis absorbida en Gy sigue la secuencia del
archivo \texttt{cargo\_mctal.m} (líneas 389--395) y se implementa en
\texttt{MCTALParser.compute\_dose\_gy()}.

\subsubsection{Paso 1: División por densidad}

La dosis en MeV/cm$^3$ se convierte a MeV/g dividiendo por la densidad del
tejido en cada voxel:

\begin{equation}
    D_{\text{MeV/g}} = \frac{D_{\text{raw}}}{\rho}
\end{equation}

Donde $\rho$ se obtiene del phantom segmentado: cada índice de tejido
(30 = tejido blando, 50 = pulmón, 80 = hueso, 90 = hígado, 100 = tumor)
tiene asociada una densidad en g/cm$^3$.

\subsubsection{Paso 2: MeV a Joules}

\begin{equation}
    D_{\text{J/g}} = D_{\text{MeV/g}} \times 1.6 \times 10^{-13}
\end{equation}

La constante de conversión es:
\begin{equation}
    1\ \text{MeV} = 1.602176634 \times 10^{-13}\ \text{J} \approx 1.6 \times 10^{-13}\ \text{J}
\end{equation}

\subsubsection{Paso 3: Integración temporal}

La dosis por partícula se escala al número total de desintegraciones
durante el tiempo de tratamiento. El factor de integración es la vida
media del $^{90}$Y:

\begin{equation}
    \tau = \frac{T_{1/2}}{\ln 2} = \frac{64.1 \times 3600}{\ln 2} \approx 332\,916\ \text{s}
\end{equation}

Este factor representa la integral desde $t=0$ hasta $t=\infty$ de la
actividad que decae exponencialmente:
\begin{equation}
    \int_0^\infty A(t)\,dt = \int_0^\infty A_0 e^{-\lambda t} dt = \frac{A_0}{\lambda} = A_0 \tau
\end{equation}

\subsubsection{Paso 4: actividad total}

Multiplicando por la actividad total administrada $A$ (en Bq):

\begin{equation}
    D_{\text{J/g}} = D_{\text{J/g/partícula}} \times \tau \times A
\end{equation}

\subsubsection{Paso 5: J/g a Gy}

\begin{equation}
    D_{\text{Gy}} = D_{\text{J/g}} \times 1000\ \frac{\text{g}}{\text{kg}}
\end{equation}

\subsubsection{Ecuación completa}

Combinando todos los pasos:

\begin{equation}
    \boxed{D_{\text{Gy}} = \frac{D_{\text{raw}}}{\rho} \times 1.6\times10^{-13} \times \tau \times A \times 1000}
\end{equation}

O equivalentemente, en la implementación Python:

\begin{equation}
    D_{\text{Gy}} = D_{\text{raw}} \cdot \frac{1}{\rho} \cdot (1.6 \times 10^{-13}) \cdot \tau \cdot A \cdot 1000
\end{equation}

\subsection{Post-procesamiento}

Antes de la conversión, se aplican filtros de calidad:

\begin{itemize}
    \item \textbf{Eliminación de errores altos:} voxeles con incertidumbre
        relativa $> 1.5$ se anulan ($D=0$). Esto ocurre en regiones con
        poca estadística (lejos de la fuente).
    \item \textbf{Eliminación de negativos:} voxeles con dosis negativa
        (ruido estadístico) se ponen a cero.
    \item \textbf{Eliminación de aire:} voxeles fuera del cuerpo (índice de
        aire) se anulan.
\end{itemize}

\subsection{Incertidumbre}

La incertidumbre estadística del tally de MCNP disminuye con
$1/\sqrt{\text{NPS}}$. Para NPS típicos de $10^7$--$10^8$, la
incertidumbre en regiones de alta dosis (tumor) suele ser $<5\%$, mientras
que en regiones de baja dosis puede superar el $100\%$.

El archivo MCTAL reporta la incertidumbre relativa $e_i$ para cada voxel
como la desviación estándar relativa ($\sigma/\bar{x}$).

% ============================================================
% MÉTODO KERNEL
% ============================================================
\section{Método Kernel: Convolución con kernel precalculado}

\subsection{Fundamento teórico}

En un medio homogéneo e infinito, la dosis en un punto $\mathbf{r}$ debida
a una distribución de fuentes $S(\mathbf{r}')$ puede expresarse como:

\begin{equation}
    D(\mathbf{r}) = \iiint S(\mathbf{r}') \, K(\mathbf{r} - \mathbf{r}') \, d^3\mathbf{r}'
\end{equation}

donde $K(\mathbf{r})$ es el \textbf{kernel de dosis}: la dosis por unidad
de actividad en un punto $\mathbf{r}$ debida a una fuente puntual en el
origen. En un medio homogéneo, $K$ depende solo de la distancia radial
$r = |\mathbf{r}|$ (simetría esférica).

La convolución se evalúa numéricamente como:

\begin{equation}
    D(\mathbf{r}) = \sum_{\mathbf{r}'} S(\mathbf{r}') \, K(\mathbf{r} - \mathbf{r}')
\end{equation}

que en forma discreta es una convolución 3D:

\begin{equation}
    D = S \ast K
\end{equation}

\subsection{Generación del kernel}

El kernel se precalcula con MCNP usando una fuente puntual de $^{90}$Y en
un medio homogéneo (agua o tejido blando, $\rho \approx 1.09$ g/cm$^3$).
El resultado es una matriz 3D (típicamente $51 \times 51 \times 51$
voxeles, aunque también se usan $65 \times 65 \times 65$) que representa
la dosis en Gy para 1 Bq de actividad en el voxel central.

En \texttt{cargo\_mctal\_kernel.m}:

\begin{align}
    \text{Kernel} &= D_3 \quad \text{(MeV/cm}^3 \text{ del tally 1)} \\
    \text{Kernel} &= \text{Kernel} \times V \quad \text{(MeV)} \\
    \text{Kernel} &= \text{Kernel} \times 1.6 \times 10^{-13} \quad \text{(J)} \\
    \text{Kernel} &= \text{Kernel} \times \tau \quad \text{(J$\cdot$s, integración temporal)} \\
    \text{Kernel} &= \text{Kernel} \times A \quad \text{(J, actividad absoluta)} \\
    \text{Kernel} &= \frac{\text{Kernel}}{m} \quad \text{(J/kg = Gy)}
\end{align}

El kernel se guarda en \texttt{kernel.mat} y se carga en Python mediante
\texttt{load\_kernel\_mat()} (módulo \texttt{dose\_kernel.py}).

\subsection{Normalización del kernel}

La propiedad fundamental del kernel es que su integral sobre todo el
espacio es igual a la dosis total depositada por 1 Bq:

\begin{equation}
    \iiint K(\mathbf{r})\,d^3\mathbf{r} = \frac{\bar{E}_\beta \cdot 1}{m}
\end{equation}

Sin embargo, en la implementación actual en Python, el kernel \textbf{no se normaliza} a suma
unitaria. Esto se debe a que \texttt{kernel.mat} ya incluye el factor $\tau \times 10^9$ (vida
media $\times$ GBq), y normalizarlo rompería la escala absoluta de dosis:

\begin{equation}
    K_{\text{uso}} = K \quad (\text{sin normalizar})
\end{equation}

La convolución con el kernel sin normalizar produce directamente dosis en Gy,
sin necesidad de multiplicar por $\tau$ posteriormente.

\subsection{Centrado del kernel}

El kernel generado por MCNP puede tener el máximo desplazado del centro
debido a la discretización del mesh tally. Es necesario centrarlo para
que la convolución FFT funcione correctamente:

\begin{equation}
    K_{\text{centrado}} = \texttt{np.roll}(K, \text{shift})
\end{equation}

donde \texttt{shift} = centro\_deseado $-$ posición\_del\_máximo.

\subsection{Convolución 3D vía FFT}

La convolución directa en el espacio real tiene costo
$\mathcal{O}(N^3 M^3)$ para una imagen de $N^3$ voxeles y kernel de
$M^3$ voxeles. Con FFT, el costo se reduce a
$\mathcal{O}(N^3 \log N)$.

\subsubsection{Algoritmo}

\begin{enumerate}
    \item \textbf{Padding simétrico:} se extiende la imagen de actividad
        reflejando los bordes (\texttt{mode='symmetric'}) para evitar
        artefactos de borde:

        \begin{equation}
            S_{\text{pad}} = \texttt{np.pad}(S, (k_r, k_r), \text{mode='symmetric'})
        \end{equation}

        donde $k_r = \lfloor M/2 \rfloor$ es el radio del kernel.

    \item \textbf{Tamaño FFT óptimo:} se calcula el tamaño óptimo para
        FFT real:

        \begin{equation}
            \text{shape}_{\text{full}} = \text{shape}(S_{\text{pad}}) + \text{shape}(K) - 1
        \end{equation}
        \begin{equation}
            \text{shape}_{\text{FFT}} = \texttt{next\_fast\_len}(\text{shape}_{\text{full}})
        \end{equation}

    \item \textbf{FFT del kernel} (se cachea para reuso):

        \begin{equation}
            K_{\text{shifted}} = \texttt{ifftshift}(K_{\text{centrado}})
        \end{equation}
        \begin{equation}
            \hat{K} = \texttt{rfftn}(K_{\text{shifted}}, s=\text{shape}_{\text{FFT}})
        \end{equation}

    \item \textbf{FFT de la actividad:}

        \begin{equation}
            \hat{S} = \texttt{rfftn}(S_{\text{pad}}, s=\text{shape}_{\text{FFT}})
        \end{equation}

    \item \textbf{Multiplicación en frecuencia:}

        \begin{equation}
            \hat{D} = \hat{S} \times \hat{K}
        \end{equation}

    \item \textbf{Transformada inversa:}

        \begin{equation}
            D_{\text{full}} = \texttt{irfftn}(\hat{D}, s=\text{shape}_{\text{FFT}})
        \end{equation}

    \item \textbf{Recorte:} se recorta el resultado para obtener el mismo
        tamaño que la imagen original (\texttt{'same'} en terminología
        MATLAB).

        \begin{align}
            D_{\text{valid}} &= D_{\text{full}}[0:\text{shape}_{\text{full}}] \\
            D_{\text{same}} &= D_{\text{valid}}[k_r : k_r + N]
        \end{align}
\end{enumerate}

\subsection{Enmascarado: hígado + tumor + peritumoral}

La dosis solo es clínicamente relevante en el hígado, tumor y una corona
peritumoral de 1 cm. Se aplica una máscara:

\begin{equation}
    D_{\text{masked}} = D \cdot \mathbb{1}_{\text{mask}}
\end{equation}

donde:

\begin{equation}
    \text{mask} = \text{hígado} \cup \text{tumor} \cup \text{peritumoral}
\end{equation}

La región peritumoral se genera mediante dilación morfológica del tumor
con un elemento estructurante elipsoidal de radio 1 cm en cada eje:

\begin{equation}
    \mathbb{1}_{\text{peritumoral}} = \text{dilate}(\text{tumor}, r) \land \lnot \text{tumor}
\end{equation}

El radio en voxeles se calcula a partir del espaciado:

\begin{equation}
    r_{\text{vox}} = \max\left(1, \left\lfloor \frac{10\ \text{mm}}{s_i} \right\rfloor\right)
\end{equation}

\subsection{Caché del kernel FFT}

Para evitar recalcular la FFT del kernel en cada ejecución, se implementa
un caché con clave basada en el shape FFT y el hash MD5 del contenido del
kernel:

\begin{equation}
    \text{clave} = (\text{shape}_{\text{FFT}}, \text{MD5}(K_{\text{bytes}}))
\end{equation}

Esto asegura que si el kernel cambia (diferente sesión de MCNP), se
recalcule automáticamente.

% ============================================================
% COMPARACIÓN
% ============================================================
\section{Comparación de Métodos}

\begin{table}[H]
\centering
\caption{Comparación entre método MCTALL y Kernel}
\label{tab:comparacion}
\begin{tabular}{p{4cm} p{5cm} p{5cm}}
\toprule
\textbf{Aspecto} & \textbf{MCTALL (Monte Carlo)} & \textbf{Kernel (Convolución)} \\
\midrule
\textbf{Base física} & Transporte de partículas & Aproximación de fuente puntual \\
\textbf{Medio} & Heterogéneo (tejidos reales) & Homogéneo (agua/tejido blando) \\
\textbf{Tiempo de simulación} & Horas--días (MCNP) & Segundos (FFT) \\
\textbf{Precisión} & Alta (estadística $\propto 1/\sqrt{N}$) & Media (aprox. medio uniforme) \\
\textbf{Incertidumbre} & Cuantificada por MCNP & No cuantificada \\
\textbf{Aplicación clínica} & Gold standard & Rápida, planificación \\
\textbf{Recursos} & Cluster/GPU & PC estándar \\
\textbf{Dependencia} & Archivo MCTAL de MCNP & Kernel precalculado + PET \\
\bottomrule
\end{tabular}
\end{table}

\subsection{¿Cuándo usar cada método?}

\begin{itemize}
    \item \textbf{MCTALL:} Cuando se necesita máxima precisión y el medio
        es heterogéneo (ej. cerca de interfaces pulmón-hígado, o cuando
        hay inhomogeneidades significativas). También cuando se quiere
        validar el método kernel.
    \item \textbf{Kernel:} En la práctica clínica diaria, donde la rapidez
        es esencial. Apropiado para $^{90}$Y porque las partículas $\beta^-$
        tienen alcance corto y el medio (hígado) es relativamente homogéneo.
\end{itemize}

% ============================================================
% MODELO MIRD
% ============================================================
\section{Anexo: Modelo MIRD (Partition Model)}

Además de los dos métodos principales, 3Dosim implementa el modelo MIRD
de partición (\emph{partition model}) para estimaciones analíticas rápidas.

\subsection{Ecuaciones MIRD}

\begin{align}
    T/N &= \frac{\text{conteo\_PET\_tumor}}{\text{conteo\_PET\_higado}} \\
    \text{FU}_{\text{normal}} &= (1 - \text{SF}) \frac{V_{\text{hígado}}}{T/N \cdot V_{\text{tumor}} + V_{\text{hígado}}} \\
    \text{FU}_{\text{tumor}} &= (1 - \text{SF}) \frac{T/N \cdot V_{\text{tumor}}}{T/N \cdot V_{\text{tumor}} + V_{\text{hígado}}} \\
    D_{\text{hígado}} &= \frac{A \cdot k \cdot \text{FU}_{\text{normal}}}{m_{\text{hígado}}} \\
    D_{\text{tumor}} &= \frac{A \cdot k \cdot \text{FU}_{\text{tumor}}}{m_{\text{tumor}}}
\end{align}

donde:
\begin{itemize}
    \item $k = 48.98$ J/s es la constante MIRD para $^{90}$Y
    \item SF = shunt pulmonar (fracción)
    \item $V$ = volumen, $m$ = masa
    \item $A$ = actividad administrada en GBq
\end{itemize}

Este modelo asume que la captación se distribuye uniformemente entre
hígado sano y tumor según la relación $T/N$ derivada del PET. Es útil
para planificación rápida pero no proporciona un mapa de dosis 3D.

% ============================================================
% REFERENCIAS
% ============================================================
\section{Referencias}

\begin{enumerate}
    \item Gulec, S.A., et al. ``Hepatic radiation dosimetry in
        $^{90}$Y-microsphere treatment.'' \textit{Med. Phys.}, 2003.
    \item Chiesa, C., et al. ``The dosimetric journey of $^{90}$Y
        radioembolization: from MIRD to voxel dosimetry.''
        \textit{EJNMMI Physics}, 2017.
    \item X.W. et al. ``$^{90}$Y microsphere dosimetry with Monte Carlo
        and convolution methods.'' \textit{Phys. Med. Biol.}, 2012.
    \item MCNP6 User's Manual, Los Alamos National Laboratory, LA-CP-13-00634.
    \item Documentación del proyecto 3Dosim:
        \texttt{cargo\_mctal.m}, \texttt{cargo\_mctal\_kernel.m},
        \texttt{mctal\_parser.py}, \texttt{dose\_kernel.py},
        \texttt{fft\_dose.py}.
\end{enumerate}

\end{document}
